\documentclass[tikz,border=6pt]{standalone}
\usepackage[no-math]{fontspec}
\usepackage[brazil]{babel}
\usepackage{amsmath,amssymb}
\IfFileExists{C:/Windows/Fonts/Lato_Regular.ttf}{%
  \setmainfont{Lato}[
    Path = C:/Windows/Fonts/,
    UprightFont = Lato_Regular.ttf,
    BoldFont = Lato_Bold.ttf,
    ItalicFont = Lato_Italic.ttf,
    BoldItalicFont = Lato_BoldItalic.ttf
  ]
  \setsansfont{Lato}[
    Path = C:/Windows/Fonts/,
    UprightFont = Lato_Regular.ttf,
    BoldFont = Lato_Bold.ttf,
    ItalicFont = Lato_Italic.ttf,
    BoldItalicFont = Lato_BoldItalic.ttf
  ]
}{%
  \RequirePackage[default]{lato}
}
\renewcommand{\familydefault}{\sfdefault}

\usepackage{xcolor}
\usepackage{tikz}
\usepackage{pgfplots}
\pgfplotsset{compat=1.18}
\usetikzlibrary{arrows.meta,positioning,shapes.geometric,fit,backgrounds,calc,decorations.pathreplacing,patterns}

% Paleta institucional UEFS / CHF662
\definecolor{uefsNavy}{HTML}{1B2A4A}
\definecolor{uefsTeal}{HTML}{0E7C7B}
\definecolor{uefsForest}{HTML}{2D6A4F}
\definecolor{uefsGold}{HTML}{C4983B}
\definecolor{uefsGoldDark}{HTML}{8F6B1E}
\definecolor{uefsAlert}{HTML}{B02A37}
\definecolor{uefsOrange}{HTML}{D95F02}
\definecolor{uefsInk}{HTML}{22252A}
\definecolor{uefsMuted}{HTML}{5C6470}
\definecolor{uefsWater}{HTML}{2E86AB}
\definecolor{uefsSand}{HTML}{D9B382}
\definecolor{bgNavy}{HTML}{F2F5F9}
\definecolor{bgTeal}{HTML}{E9F5F4}
\definecolor{bgForest}{HTML}{EEF5F0}
\definecolor{bgGold}{HTML}{FCF7EB}
\definecolor{bgAlert}{HTML}{FDF2F2}
\definecolor{bgGray}{HTML}{F8FAFC}
\definecolor{lineGray}{HTML}{CFD8DC}

\tikzset{
  >=Stealth,
  aaCard/.style={draw=#1, fill=white, rounded corners=4pt, line width=1pt, inner sep=6pt, align=center, font=\small},
  aaBox/.style={draw=#1, rounded corners=4pt, line width=1.1pt, inner sep=6pt, align=center, font=\footnotesize\sffamily, text=uefsInk},
  aaTitle/.style={font=\normalsize\bfseries\sffamily, text=uefsNavy, align=center},
  aaSub/.style={font=\footnotesize\sffamily, text=uefsMuted, align=center},
  aaArrow/.style={->, >=Stealth, line width=1.4pt},
  aaDArrow/.style={<->, >=Stealth, line width=1.4pt}
}


\begin{document}
\begin{tikzpicture}[
  panelBox/.style={
    draw=#1,
    fill=#1!7!white,
    rounded corners=6pt,
    line width=1.5pt,
    text width=6.6cm,
    minimum height=5.8cm,
    inner sep=12pt,
    align=left,
    font=\footnotesize\sffamily,
    text=uefsInk
  }
]

% Título principal (sem travessão)
\node[aaTitle, font=\large\bfseries\sffamily, text=uefsNavy] at (9.1, 3.85)
  {Como a Vegetação Estabiliza Encostas e Margens (Gray \& Sotir, 1996; Durlo \& Sutili, 2005)};

% Dois painéis bem separados
\node[panelBox=uefsTeal] (hidro) at (3.8, 0.0) {%
  \begin{minipage}[t]{6.5cm}
  \centering{\small\bfseries\color{uefsTeal}EFEITOS HIDROLÓGICOS}\\[10pt]
  \raggedright
  \textbf{1. Interceptação pela copa:}\\
  \hspace*{1em}dissipa energia cinética da chuva (evita \textit{splash}).\\[8pt]
  \textbf{2. Rugosidade e serapilheira:}\\
  \hspace*{1em}retarda o escoamento superficial e aumenta\\
  \hspace*{1em}a taxa de infiltração.\\[8pt]
  \textbf{3. Evapotranspiração:}\\
  \hspace*{1em}remove água do solo, reduz poropressão\\
  \hspace*{1em}positiva e eleva a sucção matricial.
  \end{minipage}%
};

\node[panelBox=uefsForest] (mec) at (14.4, 0.0) {%
  \begin{minipage}[t]{6.5cm}
  \centering{\small\bfseries\color{uefsForest}EFEITOS MECÂNICOS DAS RAÍZES}\\[10pt]
  \raggedright
  \textbf{1. Reforço por tração (coesão radicular $c_r$):}\\
  \hspace*{1em}raízes finas funcionam como fibras armando\\
  \hspace*{1em}a matriz do solo contra cisalhamento.\\[8pt]
  \textbf{2. Ancoragem em camadas profundas:}\\
  \hspace*{1em}raízes pivotantes atravessam o plano potencial\\
  \hspace*{1em}de ruptura e fixam o manto superficial.\\[8pt]
  \textbf{3. Retenção e contenção (arqueamento):}\\
  \hspace*{1em}troncos e raízes estruturais retêm sedimentos.
  \end{minipage}%
};

% Seta dupla ampla entre as caixas e bloco central Sinergia sem tocar bordas
\draw[aaDArrow, color=uefsGoldDark, line width=2.0pt, shorten <=8pt, shorten >=8pt]
  (hidro.east) -- (mec.west);

\node[
  draw=uefsGoldDark,
  fill=bgGold,
  rounded corners=4pt,
  line width=1.1pt,
  inner xsep=8pt,
  inner ysep=4pt,
  font=\small\bfseries\sffamily,
  text=uefsGoldDark
] at ($(hidro.east)!0.5!(mec.west) + (0, 0.58)$) {Sinergia};

\end{tikzpicture}
\end{document}
